\section{The chronological zero fibre}

Fix an element \(\rho\in\Fp^\times\) of order three.  The four-step phase is
\begin{equation}\label{eq:phase}
 \Phi(x)=2\sum_{i=0}^{3}x_i^3+
 x_0x_1+x_1x_2+x_2x_3+\rho x_3x_0.
\end{equation}
Let \(Z_p=\#\Phi^{-1}(0)\).  The exact Galois trace normalization is
\begin{equation}\label{eq:moment}
 C_{p,2}=\frac{2Z_p}{p}-2p^2,
 \qquad c_{p,2}=\frac{2C_{p,2}}{p-1}.
\end{equation}

Write \(\mathcal C=\sum x_i^3\) and
\(\mathcal Q=x_0x_1+x_1x_2+x_2x_3+\rho x_3x_0\), and denote by
\(S,R\subset\PP^3\) their zero loci and by \(X=S\cap R\).

\begin{proposition}[Direction count]\label{prop:direction}
One has
\[
 Z_p=1+\#\PP^3(\Fp)-\#S(\Fp)-\#R(\Fp)+p\#X(\Fp).
\]
\end{proposition}

\begin{proof}
For a projective direction ([v]),
\[
 \Phi(\lambda v)=\lambda^2(2\lambda\mathcal C(v)+\mathcal Q(v)).
\]
If neither homogeneous value vanishes there is one nonzero root.  If exactly
one vanishes there is none.  If both vanish all \(p-1\) nonzero scalars work.
Adding the origin proves the formula.
\end{proof}

Because \(\mu_3\subset\Fp\), the Fermat cubic surface is split; its blow-up
description gives
\[
 \#S(\Fp)=\#\PP^2(\Fp)+6p=p^2+7p+1.
\]
See also \cite[Corollary~8.3]{Massarenti2026}; since an odd prime congruent
to one modulo three is congruent to one modulo six, its split case applies.
The invertible change \(y=x_1+\rho x_3,\ w=x_1+x_3\) turns
\(\mathcal Q\) into \(x_0y+x_2w\).  Hence \(R\) is a split quadric,
\(R\simeq\PP^1\times\PP^1\), and \(\#R=(p+1)^2\).  Therefore
\begin{equation}\label{eq:zcount}
 Z_p=p^3-p^2-8p+p\#X(\Fp).
\end{equation}
